% =============================================================================
% TR Corpus Translation Benchmark — Comprehensive Technical Report v3.9
% =============================================================================
% Compiled: June 2026
% Covers: Architecture, Optimizations, Benchmarks, Issues, Codebase Evolution
% =============================================================================

\documentclass[11pt,a4paper,twoside]{article}

% ── Packages ──
\usepackage[utf8]{inputenc}
\usepackage[T1]{fontenc}
\usepackage{mathpazo}
\usepackage{geometry}
\geometry{margin=2.5cm, bindingoffset=5mm}
\usepackage{hyperref}
\usepackage{graphicx}
\usepackage{booktabs}
\usepackage{longtable}
\usepackage{enumitem}
\usepackage{fancyhdr}
\usepackage{lastpage}
\usepackage{xcolor}
\usepackage{listings}
\usepackage{tcolorbox}
\usepackage{caption}
\usepackage{float}
\usepackage{amsmath}
\usepackage{amssymb}

% ── Colors ──
\definecolor{codebg}{rgb}{0.95,0.95,0.95}
\definecolor{codeframe}{rgb}{0.7,0.7,0.7}
\definecolor{commentgreen}{rgb}{0.2,0.6,0.2}
\definecolor{criticalred}{rgb}{0.8,0.1,0.1}
\definecolor{warningorange}{rgb}{0.9,0.5,0.1}
\definecolor{infoblue}{rgb}{0.1,0.3,0.7}

% ── Code listing style ──
\lstset{
  basicstyle=\ttfamily\footnotesize,
  backgroundcolor=\color{codebg},
  frame=single,
  framerule=0.5pt,
  rulecolor=\color{codeframe},
  breaklines=true,
  numbers=left,
  numberstyle=\tiny\color{gray},
  showstringspaces=false,
  tabsize=4,
  captionpos=b,
  commentstyle=\color{commentgreen}\itshape,
}

% ── Headers & footers ──
\pagestyle{fancy}
\fancyhf{}
\fancyhead[L]{\small TR Corpus Benchmark — v3.9}
\fancyhead[R]{\small Technical Report}
\fancyfoot[C]{\thepage\ of \pageref{LastPage}}
\renewcommand{\headrulewidth}{0.4pt}

% ── Title ──
\title{%
  \huge\textbf{Turkish Corpus Translation Benchmark}\\[4pt]
  \Large\textbf{Comprehensive Technical Report — v3.9}\\[12pt]
  \normalsize Architecture, Optimization, Benchmarking, and Codebase Evolution\\
  \normalsize June 2026 --- NVIDIA H200 NVL (2$\times$141 GB HBM3e)
}
\author{%
  Benchmark Team\\
  \texttt{H200Research Project}
}
\date{\today}

\begin{document}
\maketitle

\begin{abstract}
This report documents the complete architecture, optimization stack, benchmarking
results, codebase evolution, and implementation details of the TR Corpus Translation
Benchmark --- a production-grade, model-agnostic English→Turkish translation throughput
benchmarking harness targeting NVIDIA H200 Tensor Core GPUs. The system measures
translation throughput across four backend architectures (autoregressive,
encoder-decoder), supports 16 model presets, and computes
seven quality metrics in parallel. The benchmark has undergone nine major
refinement phases totaling approximately 60,000 lines of audited, optimized, and
documented Python code. Key measured results include approximately 11,000 tok/s on NLLB-600M
encoder-decoder at batch size 64 and approximately 1,300 tok/s on TranslateGemma 4B
(BF16, pre-tokenized, PyTorch 2.6.0+cu124). Static FP8 + SmoothQuant yields
2,723 tok/s on TranslateGemma 4B with PyTorch 2.12.1+cu126 and torch.compile.
\end{abstract}

\tableofcontents
\newpage

% =============================================================================
\section{Project Overview}
% =============================================================================

\subsection{Mission Statement}

The TR Corpus Translation Benchmark answers one fundamental question: \textbf{how
many days to translate approximately 6.23 trillion English tokens into Turkish on
a dual NVIDIA H200 NVL system, at academic-grade quality?}

This benchmark is designed to be model-agnostic, platform-portable (CUDA, Apple
Silicon MPS, CPU), and production-optimized. It supports multiple translation
paradigms --- autoregressive (causal LM), encoder-decoder (NLLB, MADLAD),
diffusion-based, and vLLM engine --- through a single, unified \texttt{InferenceBackend}
protocol.

\subsection{Key Metrics}

\begin{table}[H]
\centering
\caption{Benchmark Results (Measured June 2026, Single H200 GPU)}
\label{tab:results}
\begin{tabular}{lccccr}
\toprule
\textbf{Model} & \textbf{Backend} & \textbf{Batch} & \textbf{Precision} & \textbf{TPS} & \textbf{VRAM} \\
\midrule
NLLB-200 600M     & Enc-Dec & 64    & BF16+compile & 11,000 & 1.1 GB \\
TranslateGemma 4B & AR      & 32    & BF16+pre-tok & 1,300 & 8.0 GB \\
TranslateGemma 4B & AR      & 1,749 & FP8+SmoothQuant+compile & 2,723 & 8.0 GB \\
TranslateGemma 4B & AR      & 16    & BF16+Flash SDPA & 735 & 8.0 GB \\
NLLB-200 1.3B     & Enc-Dec & 8     & BF16+Flash SDPA & 215.4 & 2.6 GB \\
NLLB-200 3.3B     & Enc-Dec & 8     & BF16+Flash SDPA & 372.5 & 6.3 GB \\
\bottomrule
\end{tabular}
\end{table}

\subsection{Scale and Scope}

\begin{itemize}[nosep]
  \item \textbf{64,393} lines of Python across 116 production files
  \item \textbf{2,560} functions and classes (73\% documented)
  \item \textbf{16} model presets (5 base architectures, multiple quantization levels)
  \item \textbf{7} quality metrics computed in parallel (3 workers)
  \item \textbf{121+} git commits refining the codebase
  \item \textbf{14} optimization features tracked via Capability Registry
  \item \textbf{25+} environment variables for fine-grained control
\end{itemize}

% =============================================================================
\section{Architecture}
% =============================================================================

\subsection{System Overview}

The benchmark follows a pipeline architecture with clearly separated subsystems:

\begin{description}[nosep, font=\ttfamily]
  \item[Entry] \texttt{benchmark/\_\_main\_\_.py} --- CLI parsing, config injection,
               run-mode dispatch
  \item[Orchestration] \texttt{orchestration/harness.py} --- owns lifecycle of all
                       subsystems: load $\rightarrow$ warmup $\rightarrow$ translate
                       $\rightarrow$ quality $\rightarrow$ report
  \item[Inference Engine] \texttt{inference/engine.py} --- model-agnostic facade,
                          backend dispatch via \texttt{ModelRegistry}
  \item[Backends] \texttt{inference/backends/} --- 14 files, 4 backend types
  \item[Data Pipeline] \texttt{data/} --- JSONL/Parquet loader, chunker, async
                       prefetch, pre-tokenization
  \item[Quality] \texttt{quality/} --- 7 parallel metrics, significance testing
  \item[Metrics] \texttt{metrics/} --- O(1) rolling throughput, GPU/system sampling
  \item[Observability] \texttt{observability/} --- Prometheus/Graphana exporter
  \item[Reporting] \texttt{reporting/} --- aggregation, extrapolation, JSON/Markdown
  \item[Quantization] \texttt{quantization/} --- SmoothQuant, QAT, static FP8
\end{description}

\subsection{The Backend Protocol}

Every backend implements the abstract \texttt{InferenceBackend} protocol
(\texttt{protocol.py}):

\begin{lstlisting}[language=Python, caption=InferenceBackend Contract]
class InferenceBackend(ABC):
    model_type: ModelType
    capabilities: ModelCapability  # IntFlag bitmask
    display_name: str

    @abstractmethod def load(self) -> None: ...
    @abstractmethod def warmup(self, batches=20) -> None: ...
    @abstractmethod def translate_batch(self, batch) -> BatchGenerationOutput: ...
    @abstractmethod def is_loaded(self) -> bool: ...
    # Optional: encode_source, score_candidates, get_token_log_probs,
    #           close, kv_cache_config
\end{lstlisting}

\subsection{Backend Dispatch}

\texttt{ModelRegistry.create\_backend} dispatches in fixed priority order:

\begin{enumerate}[nosep]
  \item Explicit override --- \texttt{extra["backend\_type"]} (supports \texttt{vllm})
  \item Auto-detect via HuggingFace model config
  \item Custom plugin --- gated behind \texttt{TR\_ALLOW\_UNTRUSTED\_PLUGINS=1}
  \item Fallback $\rightarrow$ AutoregressiveBackend
\end{enumerate}

\subsection{Backend Implementations}

\begin{table}[H]
\centering
\caption{Backend Implementations}
\label{tab:backends}
\begin{tabular}{lp{2.5cm}p{2.5cm}p{4cm}}
\toprule
\textbf{Backend} & \textbf{Model Type} & \textbf{Lines} & \textbf{Key Features} \\
\midrule
Autoregressive CUDA & AUTOREGRESSIVE & Production & Static FP8, Flash SDPA,
                     torch.compile, spec decode, PagedAttn, FA3 \\
Autoregressive MPS & AUTOREGRESSIVE & Production & Thin subclass of CUDA backend \\
NLLB CUDA & ENCODER\_DECODER & Production & torch.compile, attn\_implementation passthrough \\
NLLB MPS & ENCODER\_DECODER & Production & Direct-to-Metal loading \\
Diffusion & DIFFUSION & Experimental & T-step denoising, batched CFG, CUDA graph \\
vLLM & VLLM & Gated/Disabled & CUDA version conflict (vLLM wheels = CUDA 13.x) \\
DataParallel & AUTOREGRESSIVE & Production & N replica backends, 1 per GPU \\
Custom Plugin & CUSTOM & Gated & 	exttt{TR\_ALLOW\_UNTRUSTED\_PLUGINS=1} required \\
\bottomrule
\end{tabular}
\end{table}

\subsection{Dispatcher Pattern}

Both AR and NLLB backends use a \texttt{HardwareDispatcherBackend} (defined in
\texttt{protocol.py}) that transparently delegates to CUDA or MPS implementations
via \texttt{\_\_getattr\_\_/\_\_setattr\_\_} proxy methods. The dispatcher files
are 28 and 27 lines respectively, down from 96 and 95 --- a 71\% reduction through
deduplication.

\subsection{The Two Optimization Stacks}

\begin{description}[nosep]
  \item[Stack A --- AR Backend] Static FP8 + SmoothQuant (always on CUDA),
        torch.compile (version-gated, PT $\geq$ 2.12), Flash SDPA,
        cudaMallocAsync (when compile $\neq$ reduce-overhead), speculative
        decoder (opt-in via \texttt{---speculative}), PagedAttention (opt-in
        via \texttt{---paged-attention}). v3.7 permanently deleted 6 dead
        modules (CUDA graphs, Triton kernels, JIT CUDA/Metal, INT8 KV-cache,
        TensorRT --- 3,766 lines).
  \item[Stack B --- ContinuousBatcher + PagedAttention] The harness-owned
        CB path builds its own \texttt{PagedKVCache} and \texttt{PagedCache}
        (DynamicCache-compatibility shim). Gated behind
        \texttt{---continuous-batching ---paged-attention} (CUDA only).
        CB skips torch.compile (PagedCache non-Dynamo-traceable).
        8,192-block pool on H200.
\end{description}

% =============================================================================
\section{Optimization Stack}
% =============================================================================

\subsection{Active Optimizations (Hot Path)}

\begin{table}[H]
\centering
\caption{Active Optimizations on CUDA/H200}
\label{tab:optimizations}
\begin{tabular}{llp{3cm}p{4cm}}
\toprule
\textbf{\#} & \textbf{Feature} & \textbf{Activation} & \textbf{Impact} \\
\midrule
1  & torch.compile & PT $\geq$ 2.12, mode=default (2.12-2.13) / reduce-overhead (2.14+) & Inductor kernel fusion, frame-level graphs \\
2  & Static FP8 & Always on CUDA (TR\_SKIP\_FP8=1 to skip) & 2$\times$ memory bandwidth, dequant-on-read \\
2b & SmoothQuant & Auto on CUDA (TR\_SKIP\_SMOOTHQUANT=1 to skip) & Migrates activation outliers into weights \\
2c & Pre-tokenized Parquet & Enforced & Skips CPU tokenization, +60\% TPS \\
3  & Flash SDPA & CUDA default & 1.17--1.23$\times$ throughput \\
3b & FlashAttention-3 & \texttt{---use-flash-attention} (Hopper only) & WGMMA/TMA kernels, 1.5--2$\times$ attention \\
8  & cudaMallocAsync & Auto when compile $\neq$ reduce-overhead & Stream-ordered, zero-fragmentation \\
10 & Speculative Decode & \texttt{---speculative} (opt-in) & Layers[D:L] verify, ~25\% compute savings \\
16 & PagedAttention & \texttt{---paged-attention} (opt-in) & Block-level KV-cache, 40--70\% memory savings \\
18 & Continuous Batching & \texttt{---continuous-batching ---paged-attention} & Chunked prefill, priority scheduling \\
21 & device\_map="auto" & Model > 10\% of GPU memory & HF/accelerate pipeline parallelism \\
24 & NCCL P2P & Always on CUDA & NVLink/NVSwitch direct GPU-GPU \\
\bottomrule
\end{tabular}
\end{table}

\subsection{Static FP8 Weight Quantization}

Static FP8 replaces all \texttt{nn.Linear} layers with \texttt{StaticFP8Linear},
which stores weights in FP8 E4M3 format and dequantizes to BF16 on-chip during
the forward pass. Key properties:

\begin{itemize}[nosep]
  \item \textbf{Zero per-token overhead} --- dequantization is absorbed into the
    memory read pipeline
  \item \textbf{2$\times$ memory bandwidth} --- weights occupy half the HBM
  \item \textbf{lm\_head excluded} --- FP8 precision loss on vocabulary projection
    hurts ranking quality
  \item \textbf{TE fused kernel attempted first} --- falls back to static if
    Transformer Engine is unavailable
  \item \textbf{SmoothQuant calibration} runs automatically beforehand, smoothing
    activation outliers into weights at $\alpha = 0.50$
\end{itemize}

\subsection{SmoothQuant Calibration Pipeline}

SmoothQuant (Xiao et al., ICML 2024) migrates activation outliers into weights
before FP8 quantization:

\begin{equation}
  Y = XW = (X \cdot \text{diag}(s)^{-1}) \cdot (\text{diag}(s) \cdot W) = \hat{X} \cdot \hat{W}
\end{equation}

where $s_j = \max(|X_j|)^\alpha / \max(|W_j|)^{(1-\alpha)}$ for each channel $j$.

The calibration pipeline:
\begin{enumerate}[nosep]
  \item ActivationCapture registers forward hooks on all \texttt{nn.Linear} layers
  \item A small slice of input data (50 documents) runs through the model
  \item Activations are transferred to CPU in batch (deferred, not per-layer)
  \item SmoothQuant scales are computed and applied in-place to model weights
  \item The calibrator's \texttt{stop()} call uses \texttt{try/finally} to
    guarantee hook removal even on calibration failure
\end{enumerate}

\subsection{Speculative Decoding}

Self-speculative decoding splits the model into draft layers [0:D] and verify
layers [D:L]:

\begin{enumerate}[nosep]
  \item \textbf{Draft}: Early layers run autoregressively for K steps (cheap,
    no KV writes past layer D)
  \item \textbf{Verify}: All K candidates are processed by layers [D:L] + norm
    + lm\_head in one batched forward pass
  \item \textbf{Accept}: GPU-side vectorized comparison of draft vs. verify
    token IDs; accepted prefix until first mismatch
\end{enumerate}

Key implementation details:
\begin{itemize}[nosep]
  \item \textbf{Fixed in v3.9}: Draft KV is propagated between steps
    (\texttt{current\_kv = draft\_cache}) --- previously each step was an
    independent single-token prediction, destroying acceptance rate
  \item Verify runs \textbf{only layers [D:L]}, not the full model
  \item GPU-side comparison via \texttt{main\_preds == draft\_t} --- single
    vectorized op, no per-token \texttt{.item()} syncs
  \item Pre-allocated position and token buffers avoid per-step GPU allocations
  \item Greedy-only (temperature=0.0 --- the default for translation)
  \item Default: 8 draft / 26 verify layers for Gemma 34-layer, K=3
\end{itemize}

\subsection{PagedAttention + Continuous Batching}

vLLM-style block-level virtual memory for KV-cache management:

\begin{itemize}[nosep]
  \item \textbf{Block size}: 16 tokens $\times$ \texttt{num\_kv\_heads} $\times$
    \texttt{head\_dim} $\times$ 2 (K+V) $\times$ 2 bytes (BF16)
  \item \textbf{Block pool}: 32,768 blocks on H200 ($\approx$ 4.8 GB for 4B model)
  \item \textbf{PagedCache}: Drop-in \texttt{DynamicCache} replacement --- no model
    code changes required
  \item \textbf{Chunked prefill}: 512-token chunks, decode-priority scheduling
  \item \textbf{Heap-based priority}: $\texttt{score} = (\texttt{prompt\_len}/\texttt{max}) \times 0.5 + (\texttt{wait}/\texttt{max}) \times 0.5$
  \item \textbf{Fixed in v3.9}: Deferred heap rescore --- rebuilds only every 16
    dequeues (94\% reduction in O(N) heapify overhead)
  \item \textbf{Fixed in v3.9}: PagedKVCache.write() slice-index bug --- when
    start\_pos \% 128 $\neq$ 0, KV data was written to wrong block positions
\end{itemize}

\subsection{FlashAttention-3 on Hopper}

Wired in v3.9 with full Hopper (SM90) detection:

\begin{lstlisting}[language=Python, caption=FA3 Detection and Wiring]
if _fa3 and _is_hopper:
    torch.backends.cuda.enable_flash_attention(True)
    if hasattr(flash_attn, "flash_attn_func"):
        self._fa3_func = flash_attn.flash_attn_func
        # WGMMA async copies + TMA for KV-cache loads
        # 1.5-2x attention throughput on H200
\end{lstlisting}

% =============================================================================
\section{Quality Evaluation}
% =============================================================================

\subsection{Metrics Pipeline}

Seven quality metrics are computed in parallel via \texttt{ThreadPoolExecutor}
(3 workers, 600s per-metric timeout):

\begin{table}[H]
\centering
\caption{Quality Metrics}
\label{tab:metrics}
\begin{tabular}{lllp{3cm}}
\toprule
\textbf{Metric} & \textbf{Model} & \textbf{Type} & \textbf{Target} \\
\midrule
BERTScore & bert-base-multilingual-cased & Reference-based neural & $\geq$ 0.55 \\
COMET-22 & wmt22-comet-da & Reference-based neural & $\geq$ 0.72 \\
COMET-Kiwi & wmt22-cometkiwi-da & Reference-free neural & $\geq$ 0.60 \\
xCOMET-lite & XCOMET-XL & Reference-free neural QE & $\geq$ 0.72 \\
MetricX-24 & metricx-24-hybrid-large-v2p6 & Reference-based MQM error & $\leq$ 1.50 \\
BLEU & sacrebleu corpus\_bleu & N-gram overlap & $\geq$ 25.0 \\
chrF++ & sacrebleu corpus\_chrf (char\_order=4) & Character+word n-gram & $\geq$ 54.0 \\
\bottomrule
\end{tabular}
\end{table}

\subsection{Statistical Significance Testing}

Paired bootstrap significance testing (WMT methodology, Koehn 2004) is implemented
in \texttt{quality/significance.py}:

\begin{itemize}[nosep]
  \item $B = 10{,}000$ bootstrap resamples (default)
  \item 95\% confidence level (configurable)
  \item Paired differences $d_i = \text{score}_B^{(i)} - \text{score}_A^{(i)}$
  \item Percentile CI: lower at $B \times \alpha/2$, upper at $B \times (1-\alpha/2)$
  \item $p$-value from bootstrap distribution of absolute means
  \item Deterministic with \texttt{seed} parameter
\end{itemize}

% =============================================================================
\section{Data Pipeline}
% =============================================================================

\subsection{Pre-tokenization}

Pre-tokenization is \textbf{strictly enforced} in v3.8+. Missing pyarrow raises
\texttt{RuntimeError} --- the dynamic tokenization fallback is disabled:

\begin{lstlisting}[language=Python, caption=Pre-tokenization Enforcement]
if not HAS_PARQUET:
    raise RuntimeError(
        "Parquet support is disabled or pyarrow is missing. "
        "Install pyarrow: pip install pyarrow>=17.0.0"
    )
\end{lstlisting}

Cache invalidation uses $\text{SHA256}(\text{model\_path} + \text{tokenizer\_hash}
+ \text{max\_input\_tokens} + \text{input\_files})[:16]$.

\subsection{Async Pipeline}

\begin{itemize}[nosep]
  \item Lock-free thread-local tokenizers (up to 8 workers)
  \item Pinned-memory tensors for DMA-accelerated CPU$\rightarrow$GPU transfers on CUDA
  \item Double-buffered batch prefetch (hides 2--5ms CPU latency behind GPU compute)
  \item Prefetch worker shutdown uses queue timeouts to prevent deadlock
  \item Pre-tokenized Parquet loader as drop-in replacement for dynamic pipeline
\end{itemize}

\subsection{External Sort Shuffle}

\begin{itemize}[nosep]
  \item K-way merge for datasets exceeding 2 GiB memory budget
  \item Temp files with fsync + atomic rename for crash safety
  \item 256-file descriptor budget (configurable via \texttt{SHUFFLE\_MAX\_OPEN\_RUNS})
  \item Deterministic given seed + sorted file order
\end{itemize}

% =============================================================================
\section{Codebase Evolution — The Nine-Phase Cleanup}
% =============================================================================

\subsection{Phase 1: Dead Code \& Stale Reference Removal (v3.7--v3.8)}

\begin{itemize}[nosep]
  \item Deleted \texttt{scripts/run\_models.py} (525 lines, zero importers)
  \item Purged all TensorRT references from 7 files (Makefile, run.sh, setup.sh,
        README, Python docstrings) --- backend was deleted in v3.7 but references
        survived
  \item Removed \texttt{\_use\_int8\_kv\_cache} dead flag, \texttt{perf\_regression}
        capability label
  \item Fixed docstring lies in \texttt{speculative.py} (env var gate was removed),
        \texttt{\_\_main\_\_.py} (stale line ref), \texttt{inference/\_\_init\_\_.py}
        (wrong version/gating)
  \item Fixed references to 4 deleted docs (\texttt{ARCHITECTURAL\_FLAWS.md},
        \texttt{MEASUREMENT\_PLAN.md}, \texttt{FP8\_TE\_CUDA\_ISSUES.md},
        \texttt{PRETOKENIZATION\_PLAN.md})
  \item Updated paths for moved docs (PRD\&SDD\&SRS/, Research/)
\end{itemize}

\subsection{Phase 2: Architecture Alignment (v3.8)}

\begin{itemize}[nosep]
  \item Fixed SmoothQuant env var mismatch: \texttt{TR\_SMOOTHQUANT} $\neq$
        \texttt{TR\_SKIP\_SMOOTHQUANT} (CLI \texttt{---smoothquant} was a no-op)
  \item All 18 stale \texttt{autoregressive.py:line} references updated to
        \texttt{autoregressive\_cuda.py:line} (file split into CUDA/MPS dispatchers)
  \item \texttt{harness.py} and \texttt{pipeline.py} line references verified/corrected
  \item Feature \#20 clarified from ``\checkmark'' to ``\checkmark\ (gated)''
  \item Feature \#36 ``5 metrics'' $\rightarrow$ ``6 metrics'' + ``3 workers''
  \item Capability Registry description corrected
\end{itemize}

\subsection{Phase 3: Critical Correctness Bugs (v3.8--v3.9)}

\textbf{Seven real bugs fixed --- some with silent data corruption:}

\begin{table}[H]
\centering
\caption{Critical Bugs Fixed}
\label{tab:bugs}
\begin{tabular}{lp{5cm}p{6cm}}
\toprule
\textbf{Bug} & \textbf{File} & \textbf{Impact} \\
\midrule
SmoothQuant hook leak & \texttt{smoothquant.py:292} & Forward hooks permanently attached on calibration crash $\rightarrow$ exponential memory leak \\
FP8 near-zero weight corruption & \texttt{precision.py:375} & Weight matrices with $w_{\max} < 10^{-6}$ silently quantized to noise \\
GPU KV-cache leak & \texttt{autoregressive\_cuda.py:1487} & $\approx$33 GB dead tensors after PagedAttention prefill swap \\
PagedAttention write-index bug & \texttt{paged\_attention.py:254} & KV data written to wrong block positions when token\_pos \% 128 $\neq$ 0 \\
Speculative buffer aliasing & \texttt{speculative.py:630-647} & Position IDs and token inputs shared buffer $\rightarrow$ embedding position numbers as tokens \\
Speculative max\_new overshoot & \texttt{speculative.py:688} & Accept loop produced up to K$-$1 extra tokens beyond max\_new \\
Prefill exception swallow & \texttt{continuous\_batcher.py:651} & \texttt{\_prefill\_chunks} caught exceptions, rolled back, never re-raised $\rightarrow$ sequences decoded with incomplete KV \\
\bottomrule
\end{tabular}
\end{table}

\subsection{Phase 4: Performance Killers (v3.8--v3.9)}

\begin{itemize}[nosep]
  \item \textbf{Speculative \texttt{.item()} syncs}: GPU$\rightarrow$CPU sync on every
    accept-loop iteration. Fixed: single vectorized GPU comparison.
  \item \textbf{Speculative \texttt{torch.tensor()} allocs}: 2,000+ \texttt{cudaMalloc}
    calls per batch. Fixed: pre-allocated buffers.
  \item \textbf{ContinuousBatcher heap rebuild}: O(N) heapify on every dequeue.
    Fixed: deferred rescoring (every 16 dequeues, 94\% reduction).
  \item \textbf{ContinuousBatcher \texttt{.item()}}: Per-sequence GPU sync.
    Fixed: single \texttt{.cpu().tolist()} before the loop.
\end{itemize}

\subsection{Phase 5: Deep Architecture Bugs (v3.9)}

\begin{itemize}[nosep]
  \item \textbf{\texttt{or}-chain zero-handling}: \texttt{architecture.py} treated
    config values of 0 as absent. Fixed: \texttt{\_first\_non\_none()} helper.
  \item \textbf{Extrapolation CI scale mixing}: \texttt{rel\_uncertainty = se/mean}
    but days computed from median. Fixed: consistent denominator.
  \item \textbf{p99 percentile returns maximum}: \texttt{int(100*0.99)=99} =
    index 99/100 = max. Fixed: \texttt{ceil(n*0.99)-1}.
  \item \textbf{MPS double-call}: \texttt{\_mps\_gpu\_metrics()} called twice per
    sample. Fixed: unpack once.
\end{itemize}

\subsection{Phase 6: Documentation Overhaul (v3.9)}

\begin{itemize}[nosep]
  \item 65\% $\rightarrow$ 73\% function/class documentation coverage (306/416)
  \item 97\% coverage on critical files (speculative, paged\_attention, protocol,
    harness, etc.)
  \item Every constructor now documents parameters
  \item Every \texttt{load()}, \texttt{warmup()}, \texttt{translate()} method
    documents its contract
  \item Full module-level docstrings with architecture diagrams (speculative),
    memory savings calculations (paged\_attention), cache key computation
    (pretokenizer), statistical methodology (extrapolation)
\end{itemize}

\subsection{Phase 7: Scripts \& Config Audit}

\begin{itemize}[nosep]
  \item \textbf{CRITICAL}: \texttt{nllb\_cuda.py} NameError --- \texttt{config}
    instead of \texttt{self.config} (regression from Phase 5)
  \item \textbf{HIGH}: \texttt{evaluate.py} --- chrF/BLEU/spBLEU/Morph-BLEU scores
    overwritten per-model; every model got the last model's scores
  \item \textbf{HIGH}: \texttt{run\_nllb.py} --- encoder-decoder output token
    count subtracted encoder input length $\rightarrow$ negative TPS
  \item \textbf{MEDIUM}: subprocess timeouts hardcoded at 120--300s on 120s
    benchmark windows $\rightarrow$ 2.5--16$\times$ overrun. Fixed: dynamic
    \texttt{remaining = max(10, int(RUN\_DURATION - elapsed))}
  \item Benchmark script exception swallowing, dead nllb-660m preset,
    NLLB MPS missing \texttt{attn\_implementation}, \texttt{local\_kwargs(".")}
    fix, flash-attn-3 version pin
\end{itemize}

\subsection{Phase 8: Deep Audit --- Speculative, System Metrics, QAT, vLLM}

\begin{itemize}[nosep]
  \item \textbf{CRITICAL}: System sampler never collected metrics ---
    \texttt{start()} received \texttt{start\_time} but never assigned to
    \texttt{self.\_start\_time}. Every CPU/RAM/disk metric was \texttt{None}.
  \item \textbf{CRITICAL}: Speculative draft didn't auto-regress ---
    \texttt{current\_kv = past\_kv} set once, never updated. Each of K steps
    was an independent single-token prediction.
  \item \textbf{MEDIUM}: \texttt{qat.py} had \texttt{import math} at line 333
    but used at line 184. \texttt{ctx.save\_for\_backward} stored tensors that
    backward never reads.
  \item \textbf{MEDIUM}: vLLM \texttt{tensor\_parallel\_size=0} on CPU systems.
  \item \textbf{MEDIUM}: Diffusion \texttt{target\_len=0} edge case.
  \item Empty-texts SmoothQuant guard, batch\_logger race fix, COMET-Kiwi
    separate target constant.
\end{itemize}

\subsection{Phase 9: Configuration \& Packaging}

\begin{itemize}[nosep]
  \item \texttt{pyproject.toml}: v3.6 $\rightarrow$ v3.9, torch $\geq$ 2.4.0
    $\rightarrow$ 2.12.1, description updated
  \item \texttt{config.yaml}: speculative comment updated (autoreg fixed), DP=2
  \item \texttt{Makefile}: v3.6 $\rightarrow$ v3.9, Docker tags updated
  \item \texttt{README.md}: v3.7 $\rightarrow$ v3.9
  \item \texttt{Dockerfile}: torch 2.4.0 $\rightarrow$ 2.12.1+cu126, quantization/
    and scripts/ added to COPY, JIT/TRT removed, HEALTHCHECK fixed
  \item \texttt{launch\_llama\_server.sh}: configurable \texttt{LLAMA\_MODEL\_ROOT},
    actual GPU detection
  \item \texttt{benchmark\_models.py}: CUDA env vars added to subprocess whitelist
\end{itemize}

% =============================================================================
\section{Feature Implementation: xCOMET-lite, Significance, Vocabulary Pruning}
% =============================================================================

\subsection{xCOMET-lite (190 lines)}

Reference-free neural quality estimation using \texttt{Unbabel/XCOMET-XL}:

\begin{itemize}[nosep]
  \item Half-precision (FP16/BF16) for 2$\times$ throughput on GPU
  \item Thread-safe model caching with double-checked locking
  \item Batched inference with configurable batch size (default 32)
  \item GPU auto-detection (CUDA, MPS, CPU fallback)
  \item Score normalization (0--100 $\rightarrow$ 0--1 auto-detection)
  \item Integrated into all 7 pipelines: quality benchmark, QualityResults,
    Prometheus, markdown report, harness, \texttt{\_\_init\_\_.py} exports,
    constants
\end{itemize}

\subsection{Paired Bootstrap Significance Testing (215 lines)}

WMT-standard paired bootstrap resampling:

\begin{itemize}[nosep]
  \item $B = 10{,}000$ resamples, 95\% confidence
  \item \texttt{SignificanceResult} dataclass: mean\_diff, ci\_lower, ci\_upper,
    significant, winner, p\_value\_approx
  \item \texttt{ModelComparisonReport}: multi-model, multi-metric comparison
  \item Deterministic via \texttt{seed} parameter
\end{itemize}

\subsection{Vocabulary Pruning (290 lines)}

Bi-lingual sub-vocabulary pruning for EN$\rightarrow$TR translation:

\begin{itemize}[nosep]
  \item Tokenizes representative parallel corpus, identifies active token set
  \item Builds old$\rightarrow$new ID mapping, replaces embedding + lm\_head
    in-place
  \item GPU-side remapping via CUDA lookup table
  \item $\approx$80--85\% embedding/LM-head memory reduction (256K $\rightarrow$ 50K)
  \item 3--5$\times$ faster logit projection at batch\_size=1
  \item Dual strategy: Plan A (bilingual sub-vocabulary) + Plan B (hybrid
    full-encode + pruned-decode)
\end{itemize}

% =============================================================================
\section{Hardware Platform}
% =============================================================================

\begin{table}[H]
\centering
\caption{NVIDIA H200 NVL Specifications}
\label{tab:h200}
\begin{tabular}{ll}
\toprule
\textbf{Component} & \textbf{Specification} \\
\midrule
GPU & 2$\times$ NVIDIA H200 NVL \\
VRAM per GPU & 139.80 GB HBM3e \\
Memory Bandwidth & 4.8 TB/s per GPU \\
Interconnect & NVLink (intra-node) + InfiniBand (inter-node) \\
Tensor Cores & SM90 (Hopper), FP8 native \\
Verified PyTorch & 2.12.1+cu126 (+27\% TPS over 2.6.0) \\
Verified CUDA & 12.6 \\
CPU & 64-core \\
\bottomrule
\end{tabular}
\end{table}

% =============================================================================
\section{Known Issues and Future Work}
% =============================================================================

\subsection{Resolved Issues (All Phases)}

\begin{itemize}[nosep]
  \item 7 critical correctness bugs (silent data corruption)
  \item 4 performance bugs (.item() syncs, tensor allocs, heap rebuild, per-seq sync)
  \item 27 stale doc claims fixed
  \item 6 dead modules deleted (3,766 lines)
  \item 5 capability flags corrected/removed
  \item 3 features implemented (xCOMET-lite, significance testing, vocabulary pruning)
  \item 1,894 lines of duplicate code eliminated (autoregressive\_mps.py dedup)
  \item 191 lines of duplicate dispatcher logic eliminated
  \item All 14 optimization features verified and documented
\end{itemize}

\subsection{Documented Limitations}

\begin{itemize}[nosep]
  \item \textbf{Tensor parallelism}: \texttt{apply\_tensor\_parallelism} defined
    but never called. Multi-GPU uses HF \texttt{device\_map="auto"} or DP=2.
  \item \textbf{Tree attention}: Speculative decoding processes sequences
    one-at-a-time. Batched verification requires tree-attention kernel (future).
  \item \textbf{Determinism}: Full determinism requires \texttt{---deterministic}
    mode (not yet implemented).
  \item \textbf{Checkpoint atomicity on NFS}: \texttt{os.rename} is not atomic
    on NFS/S3 mounts. Local ext4/xfs only.
  \item \textbf{FP8 TE}: Transformer Engine is broken on pip venvs (cuBLASLt
    crash). Static FP8 + SmoothQuant is the production path.
  \item \textbf{vLLM}: Disabled due to CUDA version conflicts (vLLM wheels
    compiled for CUDA 13.x, H200 host has CUDA 12.x).
\end{itemize}

% =============================================================================
\section{Codebase Statistics}
% =============================================================================

\begin{table}[H]
\centering
\caption{Codebase Statistics (v3.9)}
\label{tab:stats}
\begin{tabular}{lr}
\toprule
\textbf{Metric} & \textbf{Value} \\
\midrule
Total Python files & 116 \\
Total lines of code & 64,393 \\
Functions and classes & 2,560 \\
Documentation coverage & 73\% (97\% on critical files) \\
Model presets & 16 (5 base + quantization variants) \\
Backend implementations & 10 (across 14 files) \\
Quality metrics & 7 \\
Prometheus metrics & 24 \\
Environment variables & 25+ \\
Test files & $\approx$27 \\
\bottomrule
\end{tabular}
\end{table}

% =============================================================================
\section*{Acknowledgments}

This report documents the collective work of nine major refinement phases on the
TR Corpus Translation Benchmark codebase. The benchmark builds upon the following
research:

\begin{itemize}[nosep]
  \item Xiao et al., ``SmoothQuant: Accurate and Efficient Post-Training
    Quantization for Large Language Models'', ICML 2024.
  \item Kwon et al., ``Efficient Memory Management for Large Language Model
    Serving with PagedAttention'', SOSP 2023.
  \item Leviathan et al., ``Fast Inference from Transformers via Speculative
    Decoding'', ICML 2023.
  \item Guerreiro et al., ``xCOMET: Transparent Machine Translation Evaluation
    through Fine-grained Error Detection'', TACL 2024.
  \item Koehn, ``Statistical Significance Tests for Machine Translation
    Evaluation'', EMNLP 2004.
  \item Kim et al., ``Vocabulary Trimming for Neural Machine Translation'',
    ACL 2019.
  \item NLLB Team, ``No Language Left Behind'', 2022.
\end{itemize}

\end{document}
